% ============================================================
% Main body — follows Appendix 1 "PLANT PRODUCTION AND
% HORTICULTURE PRACTICE" outline from the guide (23 points).
% Points 3 and 14 are conditional in the guide (research
% institute / seed production) and are marked Not Applicable
% for a commercial greenhouse tomato operation — delete those
% two subsections if your report doesn't need to mention them.
%
% Target length for the whole essay body: 10-12 typed pages
% (cover page, references, appendix and declaration do not
% count toward this).
% ============================================================

\section{Circumstances of the Establishment of the Company}
\ph{Describe when and how the greenhouse company was founded, its legal form (e.g. private enterprise, family farm, cooperative), ownership history, and the motivation behind starting tomato production under glass/plastic.}

\section{Organizational Structure of the Company}
\ph{Present the organizational chart (organogram) of the company: management, production supervisors, greenhouse staff, logistics/sales. Insert the organogram as a figure in the Appendix (uncomment the example in appendix.tex) and reference it here, e.g. \texttt{Figure~\ref{fig:organogram}}, once that figure exists.}

\subsection{Research Institute Activity}
\ph{Not applicable -- the internship took place at a commercial greenhouse operation, not a research institute or experimental station.}

\section{Human Resource Characteristics of the Company}
\ph{Describe the number of employees, qualifications, seasonal vs. permanent staff, age structure, and training practices.}

\section{Founding Capital, Fixed and Current Assets}
\ph{Give the booking value of founding capital, fixed assets (greenhouse structures, irrigation systems, vehicles) and current assets (inventories, receivables), as far as the company was willing to disclose.}

\section{Liabilities and Financing}
\ph{Describe past loans/liabilities taken on by the company, their purpose (e.g. greenhouse construction, equipment purchase), and future financing needs and possibilities.}

\section{General Description of the Plant Production Unit}
\ph{Describe the greenhouse farm as a whole: total covered area, number and type of greenhouse units, location, and how the tomato-growing unit fits into the company's broader operation.}

\section{Ecological Conditions of Production}
\begin{itemize}
    \item \textbf{Climatic conditions:} \ph{regional climate, temperature/light regime, heating and ventilation needs}
    \item \textbf{Production area:} \ph{total area under tomato cultivation, in hectares/m$^2$}
    \item \textbf{Soil properties:} \ph{soil analysis data if grown in soil}
    \item \textbf{Growing media properties:} \ph{properties of substrate used, e.g. rockwool, coco coir, perlite, if soilless culture}
\end{itemize}

\section{Analysis of Production Activity and Products}
\ph{Introduce the tomato varieties/products grown (fresh market, processing, cherry/cocktail types, etc.), volumes produced, and how production activity is analyzed by the company (yield records, quality grading).}

\section{Greenhouse Types}
\ph{Describe the greenhouse structures used: covering material (glass, polycarbonate, polyethylene film), structural type (Venlo, tunnel, etc.), and growing season(s) covered.}

\section{Description of Production Profile}
\ph{Describe the production profile of the operation: tomatoes as the main crop, any secondary crops, and how production is scheduled across the year.}

\section{Cultivated Varieties and Biological Bases}
\ph{Give a detailed description of the tomato varieties/hybrids cultivated, their biological characteristics (growth habit, ripening time, disease resistance), and seed/seedling sourcing.}

\section{Cultivation Technology}

\subsection{Soil/Substrate Cultivation}
\ph{Describe soil preparation or substrate setup prior to planting.}

\subsection{Greenhouse Preparation}
\ph{Describe preparation of the greenhouse for the crop cycle: irrigation system setup, plant support/trellising system, soil or substrate covering.}

\subsection{Nutrient Supply}
\ph{Describe fertilization program: type of fertilizer/fertigation, dosage, and timing throughout the crop cycle.}

\subsection{Planting Parameters}
\ph{Give planting time, row distance, plant spacing, and (if sown directly) sowing depth and seed rate.}

\subsection{Plant Care}
\ph{Describe ongoing plant care: pruning, training, de-leafing, irrigation scheduling, pollination (e.g. bumblebee hives).}

\subsection{Plant Protection}
\ph{Describe weed control and protection against diseases and pests (integrated pest management, biological control agents, chemical treatments used).}

\subsection{Harvesting and Storage}
\ph{Describe harvesting method and frequency, grading, packaging, and post-harvest storage/cold chain.}

\subsection{Greenhouse Culture Cultivation Technology Summary}
\ph{Summarize the overall cultivation technology cycle from planting to end-of-season crop removal.}

\subsection{Seed Production, Processing and Variety Maintenance}
\ph{Not applicable -- the internship focused on tomato production for the fresh/processing market, not sowing seed production.}

\section{Machinery Park}
\ph{List the machinery used: tractors, implements, irrigation and fertigation equipment, climate control systems (heating, ventilation, CO$_2$ dosing, screens).}

\section{Introduction of Experiments}
\ph{Describe any trials or experiments observed or assisted with during the internship (e.g. variety trials, new substrate or fertigation tests), if applicable.}

\section{Work Management}

\subsection{Organizational Structure of the Farm}
\ph{Describe how the greenhouse farm is organized day-to-day.}

\subsection{Personnel Conditions}
\ph{Describe manpower, education level of workers, and work fields/roles.}

\subsection{Organizing Working Time}
\ph{Describe working hours, rest days and shift organization, especially during peak harvest.}

\subsection{Work Organization Solutions}
\ph{Describe any specific work organization solutions used in the greenhouse (e.g. harvest crews, task rotation, quality checkpoints).}

\section{Economic Analysis}
\ph{Present yields, revenue, costs, unit cost per kg of tomatoes, profitability, efficiency indices, and a simplified cash-flow plan. Present figures in a table (see Appendix).}

\section{Marketing Directions of the Products}
\ph{Describe the marketing channels for the tomatoes (wholesale markets, supermarket chains, direct sale, export), opportunities, and conditions.}

\section{Development Goals, Objectives and Plans}
\ph{Describe the company's development goals and plans, current management problems, and possible resolutions.}

\section{Use of Subsidy}
\ph{Describe any subsidies/grants used by the company (e.g. EU agricultural subsidies) and how they are administered.}

\section{Extension Services and Advocacy}
\ph{Describe extension services, advisory bodies or advocacy groups the company uses, and their impact on farming profitability.}

\section{Developmental Aims and Plans}
\ph{Summarize the company's longer-term developmental aims and plans, and your own conclusions from the internship experience.}
